% !TeX program = xelatex
% !TeX encoding = UTF-8
\documentclass[UTF8,openany,fontset=fandol]{ctexbook}
\usepackage{amsmath,amssymb,amsthm}
\usepackage{empheq}
\usepackage{upgreek}
\usepackage{pifont}
\usepackage{graphicx}
\usepackage{subfig}
\usepackage{wrapfig}
\usepackage{enumitem}
\usepackage[perpage]{footmisc}
\usepackage{indentfirst}
\usepackage[a4paper,margin=25mm]{geometry}
\usepackage[unicode,hidelinks]{hyperref}
\ctexset{
  chapter = {name = {}, number = \arabic{chapter}, afterindent = true},
  section = {name = {}, number = \thechapter.\arabic{section}, afterindent = true},
  subsection = {name = {}, number = \thesection.\arabic{subsection}, afterindent = true}
}
\hypersetup{pdftitle={高中物理学},pdfauthor={啊波呲}}
\title{高中物理学}
\author{啊波呲}
\setcounter{secnumdepth}{2}
\setcounter{tocdepth}{2}
\setlength{\parindent}{2em}
\setlength{\parskip}{0pt}
% 插图或长公式放不下时移到下一页, 本页底部留白, 不拉伸行距.
\raggedbottom
% 增大基本行距; 行内大分数较高时, 自动保留相邻行之间的净空.
\linespread{1.25}
\setlength{\lineskip}{4pt}
\setlength{\lineskiplimit}{3pt}
% 行间公式作为正文句子的一部分处理: 矮公式少留白, 高公式和带框公式保持适度间距.
% \normalsize 会重设这些长度, 所以在正文开始时和每次 \normalsize 之后都重新设置.
\newcommand{\setdisplayskips}{%
  \setlength{\abovedisplayskip}{4pt plus 2pt minus 1pt}%
  \setlength{\belowdisplayskip}{4pt plus 2pt minus 1pt}%
  \setlength{\abovedisplayshortskip}{2pt plus 1pt minus 1pt}%
  \setlength{\belowdisplayshortskip}{3pt plus 1pt minus 1pt}%
}
\AtBeginDocument{\setdisplayskips}
\AddToHook{cmd/normalsize/after}{\setdisplayskips}
\setlength{\jot}{3pt}
\setlength{\fboxsep}{3pt}
\setlist[1]{itemsep=1pt,parsep=0pt,topsep=3pt,partopsep=0pt}
\newtheoremstyle{chinesetextbook}
 {0pt}{0pt}{\normalfont}{\parindent}{\bfseries}{}{1em}
 {\thmname{#1}\thmnumber{#2}\thmnote{ (#3)}}
\theoremstyle{chinesetextbook}
\newtheorem{example}{例}[chapter]
\renewcommand{\theexample}{\arabic{example}}
\newenvironment{solution}{\par\indent\textbf{解:} \ignorespaces}{\par}
\newenvironment{cnproof}{\par\indent\textbf{证明:} \ignorespaces}{\par}
% 并列标题之间、小标题与正文之间的间隔.
\newcommand{\bre}{\ \ \ }
% 数值后的单位, 如 $5\unit{m/s}$.
\newcommand{\unit}[1]{\ \mathrm{#1}}
\let\cleardoublepage\clearpage
